\documentclass[14pt,a4paper,oneside]{extbook}
\usepackage[utf8]{vietnam}
\usepackage{amsmath,amssymb,mathrsfs,fancyhdr,enumerate,multirow,makecell}
\usepackage{tkz-euclide,tikz-3dplot,tikz,tkz-tab}
\usetikzlibrary{arrows.meta,calc,intersections,angles,quotes,backgrounds,shapes.geometric,patterns,through,shadows,positioning}
\usepackage{yhmath}
\usepackage[top=0.7cm, bottom=0.7cm, left=1.0cm, right=1.0cm, headheight=16pt, headsep=3mm, footskip=5mm]{geometry}
\usepackage[hidelinks,unicode]{hyperref}
\usepackage{qrcode}
\usepackage[many]{tcolorbox}
\usepackage{fontawesome}

\pagestyle{fancy}
\fancyhf{}
\fancyhead[L]{\footnotesize\bfseries\color{teal!85!black}\faBook~SỔ TAY KIẾN THỨC TOÁN 4}
\fancyhead[R]{\footnotesize\bfseries\color{blue!85!black}Lớp Toán Thầy Lĩnh -- loptoanthaylinh.vn}
\fancyfoot[C]{\footnotesize\bfseries\color{gray!80!black}-- Sổ tay kiến thức: Bài 50. Biểu đồ cột --}
\renewcommand{\headrulewidth}{0.4pt}

\begin{document}

% --- BANNER TIÊU ĐỀ BÀI HỌC ---
\noindent
\begin{tcolorbox}[
  enhanced,
  colback=teal!8!white,
  colframe=teal!80!black,
  arc=3mm, boxrule=1pt,
  left=3mm, right=3mm, top=2mm, bottom=2mm,
  halign=center
]
  {\Large\bfseries\color{teal!90!black}\faBookmark\ BÀI 50: BIỂU ĐỒ CỘT}\\[1.5pt]
  {\footnotesize\itshape\color{black!75} Cấu tạo 3 thành phần, kĩ thuật đọc số liệu và phân tích so sánh trên 1 trang A4}
\end{tcolorbox}

\vspace{-2.5mm}

\noindent
\begin{minipage}[t]{0.49\textwidth}\vspace{0pt}
\begin{tcolorbox}[
  enhanced,
  colback=blue!4!white,
  colframe=blue!75!black,
  arc=2mm, boxrule=0.8pt,
  left=2mm, right=2mm, top=1.8mm, bottom=1.8mm,
  title={\footnotesize\bfseries\color{white}\faBarChart\ 1. Cấu tạo của biểu đồ cột}
]
{\footnotesize
\textbf{Biểu đồ cột} dùng các cột hình chữ nhật có cùng chiều rộng để biểu diễn số liệu:
\begin{itemize}\setlength{\itemsep}{1pt}\setlength{\leftskip}{-2.5mm}
  \item \textbf{Trục ngang (nằm ngang):} Ghi tên các đối tượng thống kê (tên lớp, môn học, ngày trong tuần...).
  \item \textbf{Trục đứng (thẳng đứng):} Ghi các vạch số lượng kèm đơn vị đo (học sinh, kg, cây...). Vạch gốc luôn là số $0$.
  \item \textbf{Các cột chữ nhật:} Có \textbf{chiều rộng bằng nhau}, cách đều nhau; \textbf{chiều cao} tương ứng với số liệu của đối tượng.
\end{itemize}
}
\end{tcolorbox}
\end{minipage}
\hfill
\begin{minipage}[t]{0.49\textwidth}\vspace{0pt}
\begin{tcolorbox}[
  enhanced,
  colback=teal!4!white,
  colframe=teal!75!black,
  arc=2mm, boxrule=0.8pt,
  left=2mm, right=2mm, top=1.8mm, bottom=1.8mm,
  title={\footnotesize\bfseries\color{white}\faSearchPlus\ 2. Kĩ thuật đọc và phân tích số liệu}
]
{\footnotesize
Khi quan sát biểu đồ cột, tiến hành 3 bước:
\begin{itemize}\setlength{\itemsep}{1pt}\setlength{\leftskip}{-2.5mm}
  \item \textbf{Đọc số liệu:} Nhìn đỉnh cột dóng ngang sang trục đứng để đọc chính xác giá trị tương ứng.
  \item \textbf{So sánh:} 
  \begin{itemize}\setlength{\itemsep}{0.5pt}\setlength{\leftskip}{-2.5mm}
    \item Cột \textbf{cao nhất} $\leftrightarrow$ Số lượng lớn nhất (nhiều nhất).
    \item Cột \textbf{thấp nhất} $\leftrightarrow$ Số lượng bé nhất (ít nhất).
  \end{itemize}
  \item \textbf{Tính toán tổng hợp:}
  \begin{itemize}\setlength{\itemsep}{0.5pt}\setlength{\leftskip}{-2.5mm}
    \item Tìm hiệu chênh lệch: Lấy số liệu cột này trừ cột kia.
    \item Tính tổng số lượng của toàn bộ biểu đồ.
  \end{itemize}
\end{itemize}
}
\end{tcolorbox}
\end{minipage}

\vspace{-2.5mm}
% --- KHỐI 3: SƠ ĐỒ TƯ DUY TIKZ ---
\noindent
\begin{tcolorbox}[
  enhanced,
  colback=orange!4!white,
  colframe=orange!80!black,
  arc=2mm, boxrule=0.8pt,
  left=2mm, right=2mm, top=1.5mm, bottom=1.5mm,
  title={\footnotesize\bfseries\color{white}\faSitemap\ SƠ ĐỒ TƯ DUY: GIẢI PHẪU CẤU TRÚC BIỂU ĐỒ CỘT}
]
\centering
\begin{tikzpicture}[scale=0.85, transform shape]
  % Trục tọa độ
  \draw[-{Stealth[length=2.5mm]}, thick] (0,0) -- (6.8,0) node[right, font=\scriptsize\bfseries] {Đối tượng (Lớp)};
  \draw[-{Stealth[length=2.5mm]}, thick] (0,0) -- (0,3.3) node[above, font=\scriptsize\bfseries] {Số cây (cây)};

  % Vạch chia trục đứng
  \foreach \y/\val in {0.7/5, 1.4/10, 2.1/15, 2.8/20} {
    \draw (0.08, \y) -- (-0.08, \y) node[left, font=\tiny\bfseries] {\val};
    \draw[dotted, gray!50] (0, \y) -- (6.2, \y);
  }
  \node[left, font=\tiny\bfseries] at (-0.08, 0) {0};

  % Các cột số liệu
  \draw[fill=blue!40, draw=blue!80!black, rounded corners=1pt] (0.8, 0) rectangle (1.8, 2.1);
  \node[above, font=\scriptsize\bfseries\color{blue!90!black}] at (1.3, 2.1) {15};
  \node[below, font=\scriptsize\bfseries] at (1.3, -0.05) {4A};

  \draw[fill=teal!45, draw=teal!80!black, rounded corners=1pt] (2.6, 0) rectangle (3.6, 2.8);
  \node[above, font=\scriptsize\bfseries\color{teal!90!black}] at (3.1, 2.8) {20};
  \node[below, font=\scriptsize\bfseries] at (3.1, -0.05) {4B};

  \draw[fill=orange!45, draw=orange!80!black, rounded corners=1pt] (4.4, 0) rectangle (5.4, 1.4);
  \node[above, font=\scriptsize\bfseries\color{orange!90!black}] at (4.9, 1.4) {10};
  \node[below, font=\scriptsize\bfseries] at (4.9, -0.05) {4C};

  % Callout giải thích cấu tạo
  \node[draw=purple!80!black, fill=purple!8, rounded corners=2pt, font=\tiny\bfseries, align=center] (call1) at (6.8, 2.7) {CỘT CAO NHẤT\\(4B nhiều nhất)};
  \draw[-{Stealth}, purple!80!black, thick] (call1.west) -- (3.65, 2.7);

  \node[draw=teal!80!black, fill=teal!8, rounded corners=2pt, font=\tiny\bfseries, align=center] (call2) at (6.8, 1.2) {CỘT THẤP NHẤT\\(4C ít nhất)};
  \draw[-{Stealth}, teal!80!black, thick] (call2.west) -- (5.45, 1.2);
\end{tikzpicture}
\end{tcolorbox}

\vspace{-2.5mm}
% --- KHỐI 4: VÍ DỤ & MẸO TRÁNH BẪY ---
\noindent
\begin{minipage}[t]{0.49\textwidth}\vspace{0pt}
\begin{tcolorbox}[
  enhanced,
  colback=green!4!white,
  colframe=green!60!black,
  arc=2mm, boxrule=0.8pt,
  left=2mm, right=2mm, top=1.8mm, bottom=1.8mm,
  title={\footnotesize\bfseries\color{white}\faCheckCircle\ Ví dụ thực hành chuẩn mực}
]
{\footnotesize
\textbf{Ví dụ:} Quan sát biểu đồ trên và trả lời:
\begin{itemize}\setlength{\itemsep}{1pt}\setlength{\leftskip}{-2.5mm}
  \item Lớp 4B trồng được bao nhiêu cây?\\
  $\to$ Nhìn đỉnh cột 4B dóng sang: \textbf{20 cây}.
  \item Lớp nào trồng nhiều cây nhất?\\
  $\to$ Cột 4B cao nhất: \textbf{Lớp 4B}.
  \item Cả 3 lớp trồng được tất cả bao nhiêu cây?\\
  $\to$ Tính tổng: $15 + 20 + 10 = \mathbf{45}$ (cây).
  \item Lớp 4B trồng nhiều hơn 4C bao nhiêu cây?\\
  $\to$ Tính hiệu: $20 - 10 = \mathbf{10}$ (cây).
\end{itemize}
}
\end{tcolorbox}
\end{minipage}
\hfill
\begin{minipage}[t]{0.49\textwidth}\vspace{0pt}
\begin{tcolorbox}[
  enhanced,
  colback=red!4!white,
  colframe=red!70!black,
  arc=2mm, boxrule=0.8pt,
  left=2mm, right=2mm, top=1.8mm, bottom=1.8mm,
  title={\footnotesize\bfseries\color{white}\faExclamationTriangle\ Mẹo tránh bẫy trắc nghiệm}
]
{\footnotesize
\begin{itemize}\setlength{\itemsep}{1.5pt}\setlength{\leftskip}{-2.5mm}
  \item \textbf{Bẫy vạch chia trên trục đứng:}\\
  Phải kiểm tra xem mỗi vạch chia tương ứng với bao nhiêu đơn vị ($1, 2, 5$ hay $10$ đơn vị). Không được tự ý đếm số vạch!
  \item \textbf{Bẫy đọc nhầm tên đối tượng:}\\
  Đối chiếu cẩn thận tên đối tượng dưới chân cột với đỉnh cột trước khi đọc số liệu.
  \item \textbf{Bẫy quên đơn vị đo:}\\
  Luôn chú ý đơn vị đo ghi ở đầu trục đứng (tấn, tạ, kg, nghìn đồng...).
\end{itemize}
}
\end{tcolorbox}
\end{minipage}

\end{document}
